% ============================================================================
% LANGENTWURF LE-08a: Viajar (Reisen)
% Spanisch Klasse 7 - Sequenz 8: Vacaciones
% ============================================================================
% Tier: 1 (vollständig)
% Doppelstunde: 1-2 von 8
% Fachfarbe: #c41e3a (Karminrot)
% ============================================================================

\documentclass[11pt,a4paper]{article}

\usepackage[utf8]{inputenc}
\usepackage[T1]{fontenc}
\usepackage[ngerman]{babel}
\usepackage{geometry}
\usepackage{fancyhdr}
\usepackage{lastpage}
\usepackage{xcolor}
\usepackage{tcolorbox}
\usepackage{enumitem}
\usepackage{tabularx}
\usepackage{longtable}
\usepackage{array}
\usepackage{booktabs}
\usepackage{amssymb}

% ============================================================================
% KONFIGURATION
% ============================================================================

\newcommand{\fach}{Spanisch}
\newcommand{\fachfarbe}{c41e3a}
\newcommand{\jahrgangsstufe}{7}
\newcommand{\schulart}{Gesamtschule}
\newcommand{\stundenthema}{Viajar -- Reisen und Verkehrsmittel}
\newcommand{\reihenthema}{Vacaciones}
\newcommand{\stundennummer}{8a}
\newcommand{\reihenstunden}{8}
\newcommand{\gerniveau}{A1}
\newcommand{\lehrwerk}{Qué pasa 1}
\newcommand{\lehrwerkseite}{S. 124-128}
\newcommand{\lerngruppengroesse}{24}

% ============================================================================
% LAYOUT
% ============================================================================
\geometry{left=2.5cm, right=2.5cm, top=2.5cm, bottom=2.5cm}
\definecolor{fach}{HTML}{\fachfarbe}
\definecolor{gray}{HTML}{666666}
\definecolor{lightgray}{HTML}{f5f5f5}
\definecolor{successgreen}{HTML}{2a7a4b}
\definecolor{warningorange}{HTML}{b35c00}
\definecolor{basis}{HTML}{2a7a4b}
\definecolor{standard}{HTML}{2c5aa0}
\definecolor{erweiterung}{HTML}{7b1fa2}

\tcbuselibrary{skins,breakable}

\newtcolorbox{infobox}[1][]{
    colback=lightgray,
    colframe=fach,
    fonttitle=\bfseries,
    title=#1,
    breakable
}

\newtcolorbox{scorebox}{
    colback=white,
    colframe=fach,
    fonttitle=\bfseries,
    title=Didaktik-Score,
    breakable
}

\pagestyle{fancy}
\fancyhf{}
\fancyhead[L]{\small\textcolor{gray}{\fach{} -- Jg. \jahrgangsstufe{} -- \stundenthema}}
\fancyhead[R]{\small\textcolor{gray}{LE-\stundennummer{} | Tier 1}}
\fancyfoot[C]{\small\thepage{} / \pageref{LastPage}}
\renewcommand{\headrulewidth}{0.4pt}

% Differenzierungsmarker
\newcommand{\B}{\textcolor{basis}{\textbf{[B]}}}
\newcommand{\Std}{\textcolor{standard}{\textbf{[S]}}}
\newcommand{\E}{\textcolor{erweiterung}{\textbf{[E]}}}

% ============================================================================
% DOKUMENT
% ============================================================================
\begin{document}

% ============================================================================
% DECKBLATT
% ============================================================================
\begin{titlepage}
    \centering
    \vspace*{2cm}
    {\large\textcolor{gray}{\schulart}}\\[2cm]
    {\Huge\textcolor{fach}{\textbf{Langentwurf}}}\\[0.5cm]
    {\large LE-\stundennummer{} | Tier 1 (vollständig)}\\[2cm]
    {\LARGE\textbf{\stundenthema}}\\[0.5cm]
    {\large Sequenz: \reihenthema}\\[2cm]
    \begin{tabular}{rl}
        \textbf{Fach:} & \fach{} (\gerniveau) \\[0.3cm]
        \textbf{Klasse:} & \jahrgangsstufe \\[0.3cm]
        \textbf{Lehrwerk:} & \lehrwerk, \lehrwerkseite \\[0.3cm]
        \textbf{Doppelstunde:} & \stundennummer{} (Std. 1-2 von 8) \\[0.3cm]
    \end{tabular}
    \vfill
\end{titlepage}

\tableofcontents
\newpage

% ============================================================================
% 1. BEDINGUNGSANALYSE
% ============================================================================
\section{Bedingungsanalyse}

\subsection{Lerngruppe}
\begin{tabular}{ll}
    Kursgröße: & \lerngruppengroesse{} SuS \\
    Jahrgangsstufe: & \jahrgangsstufe \\
    Sprachniveau: & \gerniveau{} (GER) -- Anfänger \\
    Fremdsprache: & 2. Fremdsprache (nach Englisch) \\
\end{tabular}

\subsection{Vorwissen aus Sequenz 1-7}
\begin{itemize}
    \item Grundwortschatz: Begrüßung, Familie, Haus, Schule, Stadt, Einkaufen
    \item Verben: ser, estar, tener, gustar, verbos regulares (-ar, -er, -ir)
    \item Strukturen: Verneinung, Fragen, Possessivpronomen
    \item ir + a + Infinitiv (nahe Zukunft) -- wurde in Seq. 5 eingeführt
\end{itemize}

\subsection{Differenzierung (intern)}
\begin{tabular}{|l|p{3.5cm}|p{3.5cm}|p{3.5cm}|}
\hline
& \B{} \textbf{Basis} & \Std{} \textbf{Standard} & \E{} \textbf{Erweiterung} \\
\hline
\textbf{Ziel} & BR: Rezeption + Reproduktion & MR: Produktion mit Hilfe & GY: Freie Produktion \\
\hline
\textbf{Vokabeln} & 8 Verkehrsmittel nennen & + Orte zuordnen & + Aktivitäten beschreiben \\
\hline
\textbf{ir a + Inf.} & Lückentext mit Vorgabe & Sätze selbst bilden & Reiseplan schreiben \\
\hline
\end{tabular}

\vspace{0.5em}
\textbf{WICHTIG:} Diese Marker sind NUR für die Lehrkraft!

% ============================================================================
% 2. SACHANALYSE
% ============================================================================
\section{Sachanalyse}

\subsection{Thematischer Wortschatz}

\textbf{Verkehrsmittel (Los medios de transporte):}
\begin{center}
\begin{tabular}{|l|l|l|}
\hline
\textbf{Spanisch} & \textbf{Deutsch} & \textbf{Beispiel} \\
\hline
el avión & das Flugzeug & Vamos en avión a España. \\
el tren & der Zug & El tren sale a las ocho. \\
el coche & das Auto & Viajamos en coche. \\
el autobús & der Bus & Tomo el autobús al centro. \\
el barco & das Schiff & El barco va a las islas. \\
\hline
\end{tabular}
\end{center}

\textbf{Urlaubsorte (Los destinos):}
\begin{center}
\begin{tabular}{|l|l|}
\hline
\textbf{Spanisch} & \textbf{Deutsch} \\
\hline
la playa & der Strand \\
la montaña & die Berge \\
el hotel & das Hotel \\
el camping & der Campingplatz \\
\hline
\end{tabular}
\end{center}

\textbf{Urlaubsaktivitäten (Las actividades):}
\begin{center}
\begin{tabular}{|l|l|}
\hline
\textbf{Spanisch} & \textbf{Deutsch} \\
\hline
nadar & schwimmen \\
tomar el sol & sich sonnen \\
hacer senderismo & wandern \\
\hline
\end{tabular}
\end{center}

\subsection{Grammatik: ir a + Infinitiv (Wiederholung)}

\textbf{Bildung:}
\begin{center}
\begin{tabular}{|l|l|l|}
\hline
\textbf{Person} & \textbf{ir} & \textbf{+ a + Infinitiv} \\
\hline
yo & voy & a nadar \\
tú & vas & a viajar \\
él/ella/usted & va & a hacer senderismo \\
nosotros/-as & vamos & a tomar el sol \\
vosotros/-as & vais & a ir a la playa \\
ellos/ellas/ustedes & van & a quedarse en el hotel \\
\hline
\end{tabular}
\end{center}

\textbf{Verwendung:} Nahe Zukunft / geplante Handlungen
\begin{itemize}
    \item \textit{Voy a nadar.} = Ich werde schwimmen. / Ich gehe schwimmen.
    \item \textit{Vamos a viajar en avión.} = Wir werden mit dem Flugzeug reisen.
\end{itemize}

\subsection{Kulturelle Aspekte}
\begin{itemize}
    \item Spanien als beliebtes Urlaubsziel (Balearen, Kanaren, Costa Brava)
    \item Typisch spanische Urlaubsgewohnheiten (siesta am Strand, tapas am Abend)
    \item Lateinamerika: Karibik, Mexiko als Reiseziele
\end{itemize}

% ============================================================================
% 3. DIDAKTISCHE ANALYSE
% ============================================================================
\section{Didaktische Analyse}

\subsection{Stellung in der Sequenz}
\begin{tabular}{|c|l|l|}
\hline
\textbf{Block} & \textbf{Thema} & \textbf{Status} \\
\hline
\textbf{8a} & \textbf{Viajar -- Verkehrsmittel + ir a + Inf.} & \textbf{diese Stunde} \\
8b & Las vacaciones -- Urlaubsaktivitäten & folgt \\
8c & El tiempo -- Wetter + Jahreszeiten & folgt \\
8d & Vacaciones -- Sequenztest & folgt \\
\hline
\end{tabular}

\subsection{Kompetenzziele (Rahmenplan MV)}
\begin{itemize}
    \item \textbf{Kommunikativ:} Sprechen (zusammenhängendes Sprechen), Hörverstehen
    \item \textbf{Sprachlich:} Wortschatz (Reisen), Grammatik (ir a + Infinitiv)
    \item \textbf{Interkulturell:} Urlaubsziele in Spanien und Lateinamerika
    \item \textbf{Methodisch:} Wortfeldarbeit, Bildgestützte Semantisierung
\end{itemize}

\subsection{Legitimation}
\textbf{Rahmenplan Spanisch MV (Kl. 7-10):}
\begin{itemize}
    \item Kompetenzbereich: Kommunikative Kompetenz -- Sprechen
    \item Standard: ``Die SuS können über Pläne und Vorhaben sprechen und einfache Informationen zu Reisen und Urlaub austauschen.''
\end{itemize}

% ============================================================================
% 4. STUNDENZIELE
% ============================================================================
\section{Stundenziele}

\subsection{Grobziel}
Die SuS erweitern ihre \textbf{kommunikative Kompetenz im Bereich Sprechen}, indem sie Reisevokabular (Verkehrsmittel, Orte, Aktivitäten) anwenden und mit \textit{ir a + Infinitiv} über Urlaubspläne sprechen.

\subsection{Feinziele (operationalisiert)}
\begin{tabular}{|c|p{7.5cm}|c|c|}
\hline
\textbf{Nr.} & \textbf{Die SuS können...} & \textbf{AFB} & \textbf{Diff.} \\
\hline
FZ1 & mindestens 5 Verkehrsmittel auf Spanisch \textbf{benennen} und korrekt aussprechen & I & alle \\
\hline
FZ2 & Urlaubsorte und -aktivitäten den passenden Bildern \textbf{zuordnen} & I & alle \\
\hline
FZ3 & mit \textit{ir a + Infinitiv} Sätze über Urlaubspläne \textbf{bilden} & II & alle \\
\hline
FZ4 & einen einfachen Dialog über Urlaubspläne \textbf{führen} & II & \Std{}/\E{} \\
\hline
FZ5 & einen kurzen Reiseplan \textbf{verfassen} und kulturelle Aspekte \textbf{einbeziehen} & III & \E{} \\
\hline
\end{tabular}

\vspace{1em}
\textbf{AFB-Verteilung:} AFB I: 30\% | AFB II: 50\% | AFB III: 20\%

% ============================================================================
% 5. METHODISCHE ANALYSE
% ============================================================================
\section{Methodische Analyse}

\subsection{Einstieg}
\begin{itemize}
    \item \textbf{Methode:} Bildimpuls -- Collage mit Urlaubsfotos (Strand, Berge, Verkehrsmittel)
    \item \textbf{Begründung:} Aktivierung von Vorwissen, Motivation, Themenankündigung
    \item \textbf{Alternative:} Kurzvideo über spanische Urlaubsziele
\end{itemize}

\subsection{Erarbeitung}
\begin{itemize}
    \item \textbf{Methode:} Bildgestützte Semantisierung -- Vokabeln über Bilder einführen
    \item \textbf{Begründung:} Visueller Zugang erleichtert Merkfähigkeit
    \item \textbf{Scaffolding:} Wortschatzkarten mit Bild + Wort, chorisches Sprechen
\end{itemize}

\subsection{Übungsphasen}
\begin{itemize}
    \item \textbf{Methode:} Memory-Spiel (Bild-Wort) $\rightarrow$ Partnerarbeit $\rightarrow$ Dialog
    \item \textbf{Progression:} Rezeptiv $\rightarrow$ Reproduktiv $\rightarrow$ Produktiv
    \item \textbf{Differenzierung:} Wortschatzhilfen für \B{}, freie Produktion für \E{}
\end{itemize}

\subsection{Medien}
\begin{itemize}
    \item Tafel/Whiteboard: Vokabeltabelle, Grammatikkasten
    \item Lehrwerk: \lehrwerk, \lehrwerkseite
    \item Arbeitsblatt: AB-08a.pdf
    \item Lernpfad: LP-08a.html (Tablets, optional)
    \item Bildkarten: Verkehrsmittel, Orte, Aktivitäten
\end{itemize}

% ============================================================================
% 6. VERLAUFSPLANUNG
% ============================================================================
\section{Verlaufsplanung}

\textbf{1. Stunde (45 Min.)}

\begin{longtable}{|p{1cm}|p{1.5cm}|p{4cm}|p{3cm}|p{1.5cm}|p{2cm}|}
\hline
\textbf{Zeit} & \textbf{Phase} & \textbf{Lehrerhandeln} & \textbf{Schülerhandeln} & \textbf{SF} & \textbf{Material} \\
\hline
\endhead

0--5 & Einstieg & Bildcollage zeigen: ``¿Qué veis?'' Vermutungen sammeln & Bilder beschreiben, Vorwissen aktivieren & Plenum & Bildcollage \\
\hline

5--15 & Input 1 & Verkehrsmittel einführen: Bilder zeigen, Wörter nennen & Nachsprechen, Aussprache üben & Plenum & Bildkarten \\
\hline

15--25 & Übung 1 & Memory-Spiel anleiten: Bild-Wort-Zuordnung & Memory spielen, Vokabeln festigen & PA & Memory-Karten \\
\hline

25--35 & Input 2 & Urlaubsorte + Aktivitäten einführen & Vokabeln notieren, nachsprechen & Plenum & Tafel, AB \\
\hline

35--45 & Übung 2 & AB-08a Aufgabe 1+2: Zuordnung & Einzelarbeit, dann Vergleich & EA/PA & AB-08a \\
\hline
\end{longtable}

\vspace{1em}

\textbf{2. Stunde (45 Min.)}

\begin{longtable}{|p{1cm}|p{1.5cm}|p{4cm}|p{3cm}|p{1.5cm}|p{2cm}|}
\hline
\textbf{Zeit} & \textbf{Phase} & \textbf{Lehrerhandeln} & \textbf{Schülerhandeln} & \textbf{SF} & \textbf{Material} \\
\hline
\endhead

0--5 & Aktivier. & Schnelles Quiz: Verkehrsmittel benennen & Antworten & Plenum & Bildkarten \\
\hline

5--15 & Input 3 & ir a + Infinitiv wiederholen: Tabelle, Beispiele & Tabelle ergänzen, Beispiele bilden & Plenum & Tafel, AB \\
\hline

15--25 & Übung 3 & AB-08a Aufgabe 3: Sätze mit ir a + Inf. & Sätze schreiben & EA & AB-08a \\
\hline

25--35 & Anwendung & PA: Minidialog über Urlaubspläne & Dialog üben, präsentieren & PA & Karten \\
\hline

35--40 & Erweiter. & \E{}: Reiseplan schreiben (5 Sätze) & Texte schreiben & EA & AB-08a \\
\hline

40--45 & Sicherung & Merkkasten zusammenfassen, HA & Notieren & Plenum & AB-08a \\
\hline
\end{longtable}

\textbf{Legende:} EA = Einzelarbeit, PA = Partnerarbeit, GA = Gruppenarbeit

% ============================================================================
% 7. ERWARTETE SCHWIERIGKEITEN
% ============================================================================
\section{Erwartete Schwierigkeiten}

\begin{tabular}{|p{4.5cm}|p{8cm}|}
\hline
\textbf{Schwierigkeit} & \textbf{Reaktion} \\
\hline
Aussprache \textit{avión} [a'bjon] & Vorsprechen, chorisch wiederholen, Silben klatschen \\
\hline
Verwechslung \textit{voy/vas/va} & Kontrastive Übung: Personalpronomen deutlich machen \\
\hline
Fehlende Präposition \textit{en} & Regel verdeutlichen: en + Verkehrsmittel (en avión, en tren) \\
\hline
Infinitiv ohne \textit{a} & Merkspruch: ``ir A viajar'' -- das A gehört dazu! \\
\hline
Zeitproblem & Puffer: Erweiterungsaufgabe als HA aufgeben \\
\hline
\end{tabular}

% ============================================================================
% 8. HAUSAUFGABE
% ============================================================================
\section{Hausaufgabe}

\begin{itemize}
    \item Lehrwerk S. 126, Übung 4 (Vokabeln)
    \item Vokabeln lernen: Verkehrsmittel, Orte, Aktivitäten
    \item Optional: Lernpfad LP-08a.html bearbeiten
    \item \E{}: Reiseplan erweitern (10 Sätze)
\end{itemize}

% ============================================================================
% 9. LÖSUNGEN
% ============================================================================
\section{Lösungen}

\subsection{AB-08a Lösungen}

\textbf{Aufgabe 1: Verkehrsmittel benennen}
\begin{itemize}
    \item el avión, el tren, el coche, el autobús, el barco
\end{itemize}

\textbf{Aufgabe 2: Zuordnung Orte + Aktivitäten}
\begin{enumerate}[label=\alph*)]
    \item la playa $\rightarrow$ nadar, tomar el sol
    \item la montaña $\rightarrow$ hacer senderismo
    \item el hotel $\rightarrow$ descansar
    \item el camping $\rightarrow$ hacer una barbacoa
\end{enumerate}

\textbf{Aufgabe 3: ir a + Infinitiv}
\begin{enumerate}[label=\alph*)]
    \item Yo \textbf{voy a nadar} en la playa.
    \item Nosotros \textbf{vamos a viajar} en avión.
    \item Tú \textbf{vas a hacer} senderismo.
    \item Ellos \textbf{van a tomar} el sol.
\end{enumerate}

\textbf{Aufgabe 4: Dialog (Beispiel)}
\begin{itemize}
    \item A: ¿Adónde vas a ir de vacaciones?
    \item B: Voy a ir a la playa.
    \item A: ¿Cómo vas a viajar?
    \item B: Voy a viajar en avión.
    \item A: ¿Qué vas a hacer?
    \item B: Voy a nadar y a tomar el sol.
\end{itemize}

\textbf{Aufgabe 5: Reiseplan (Beispiel)}
\begin{itemize}
    \item Este verano voy a ir a España.
    \item Voy a viajar en avión.
    \item Voy a quedarme en un hotel.
    \item Voy a ir a la playa todos los días.
    \item Voy a nadar y a tomar el sol.
\end{itemize}

% ============================================================================
% 10. ANLAGEN
% ============================================================================
\section{Anlagen}

\begin{enumerate}
    \item Arbeitsblatt (AB-08a.pdf)
    \item Musterlösung (ML-08a.pdf)
    \item Lehrerhinweise (LH-08a.pdf)
    \item Lernpfad (LP-08a.html)
    \item Bildkarten Verkehrsmittel (Kopiervorlage)
    \item Memory-Karten (Kopiervorlage)
\end{enumerate}

% ============================================================================
% 11. DIDAKTIK-SCORE
% ============================================================================
\newpage
\section{Didaktik-Score}

\begin{scorebox}
\textbf{Bewertung der didaktischen Qualität nach 10 Dimensionen}\\
\textbf{Ziel-Score: $\geq$ 9.0 (Premium-Qualität)}

\vspace{1em}
\begin{tabular}{|c|l|c|c|p{4cm}|}
\hline
\textbf{\#} & \textbf{Dimension} & \textbf{Gew.} & \textbf{Score} & \textbf{Notiz} \\
\hline
1 & Differenzierung & 15\% & 9/10 & [B]/[S]/[E] qualitativ unterschieden \\
\hline
2 & Taxonomie (AFB) & 10\% & 9/10 & 30/50/20 -- ausgewogen \\
\hline
3 & Lernziel-Alignment & 10\% & 10/10 & RP MV explizit, Reisevokabular \\
\hline
4 & Aktivierung & 10\% & 9/10 & Bilder, Memory, Dialog \\
\hline
5 & Scaffolding & 10\% & 9/10 & Gestufte Hilfen, Wortschatzkarten \\
\hline
6 & Feedback-Qualität & 10\% & 9/10 & LP: elaboriertes Feedback \\
\hline
7 & Sprachliche Klarheit & 10\% & 10/10 & Kl. 7 angemessen \\
\hline
8 & Motivation & 10\% & 10/10 & Urlaubsthema -- hohe Relevanz! \\
\hline
9 & Inklusion (UDL) & 10\% & 9/10 & Bilder, Text, Spiele \\
\hline
10 & Praktikabilität & 5\% & 9/10 & 90 Min. realistisch \\
\hline
\multicolumn{3}{|r|}{\textbf{GESAMT:}} & \textbf{9.3/10} & \\
\hline
\end{tabular}

\vspace{1em}
$\checkmark$ \textcolor{successgreen}{\textbf{FREIGABE} (Score $\geq$ 9.0)}
\end{scorebox}

\end{document}
